\section{{\ProjectName}}


\begin{figure}[h]
\begin{center}
\includegraphics[width=\textwidth]{figures/main_framework_1.pdf} 
\end{center}
\vspace{-10pt}
\caption{Overview of {\ProjectName}, illustrating its unified scientific discovery pipeline organized around the Generation, Verification, and Evolution subsystems. The system operates through foundational capabilities for deep research, solution refinement, and long horizon memory, which together enable sustained autonomous scientific discovery.}
\label{img:framework}
\end{figure}
%%%%%%%%%%%%%%%%%%%%%%%%%%%%%%%%%%%%%%%%%%%%%%%%%%%%%%%%%%%%%%%%%%%%%%%%%%%%%%%%
\subsection{System Overview}

In this section, we present the system overview of {\ProjectName} as illustrated in Fig.~\ref{img:framework}. The system automates the full research cycle by integrating hypothesis formulation, methodological evaluation, and evidence driven refinement into a unified and continuously improving process. Its operation relies on foundational capabilities that support deep research, solution refinement, and long horizon memory. These capabilities are realized through agent driven reasoning and system level infrastructure and allow the system to maintain contextual continuity across iterations. With this capability structure, {\ProjectName} coordinates multiple subsystems to support autonomous, scalable, and sustained scientific discovery.

\subsubsection{Architecture}
The architecture of {\ProjectName} is organized into three core subsystems, namely the \textbf{Generation}, the \textbf{Verification}, and the \textbf{Evolution}. These subsystems form an integrated and iterative workflow. The Generation subsystem formulates hypotheses and methodological plans, the Verification subsystem evaluates these plans through computational or empirical procedures, and the Evolution subsystem incorporates the resulting evidence to update internal knowledge, strategies, and long term memory. This organization maintains a coherent flow of information and enables multi cycle autonomous operation.

\begin{itemize}
    \item \textbf{Generation:} The Generation subsystem constructs the conceptual foundation of scientific inquiry. It follows the generation and reflection paradigm of InternAgent 1.0~\citep{team2025novelseek} and is driven by the foundational capability of deep research. It conducts large scale literature analysis, scientific reasoning, and contextual integration and may invoke scientific tools when processing domain specific data. It produces structured hypotheses and methodological plans and records key reasoning traces for subsequent processing.

    \item \textbf{Verification:} The Verification subsystem evaluates the hypotheses and methodological plans produced by the Generation subsystem. Its operation is driven by the foundational capability of solution refinement, which structures the iterative search for improved procedures. It performs computational analyses, simulations, and laboratory style assessments as needed and uses historical information to guide evaluation choices. It supports parallel assessment of methodological variants and generates structured evidence for downstream refinement.

    \item \textbf{Evolution:} The Evolution subsystem refines system understanding and long term strategies based on outcomes from the Generation and Verification subsystem. It is driven by the foundational capability of memory and unifies analytical feedback with persistent knowledge management. It interprets verification results, identifies methodological limitations, updates procedural, episodic, and semantic information, and produces refined priors that guide subsequent cycles of the Generation and Verification subsystem.
\end{itemize}

These subsystems rely on the foundational capabilities introduced above in order to function coherently across extended discovery horizons. The next section presents these capabilities in detail and describes the technical methods that implement them.


\subsubsection{Foundational Capabilities}

The operation of {\ProjectName} relies on a set of foundational capabilities that allow the \textbf{Generation}, \textbf{Verification}, and \textbf{Evolution} subsystems to function coherently across extended discovery cycles. These capabilities support literature based hypothesis construction, methodological evaluation, iterative refinement, and long horizon continuity. They are implemented through the technical methods introduced in Sections \ref{sec:dr_method} to \ref{sec:memory_method} and provide the requirements for end-to-end scientific discovery.

\textbf{\textit{The first capability is deep research.}} It supports the Generation subsystem by enabling large scale retrieval, integration, and structuring of cross disciplinary scientific knowledge. Section \ref{sec:dr_method} introduces the search mechanisms and structured representations that realize this capability.

\textbf{\textit{The second capability is solution refinement.}} It supports the Verification subsystem by guiding the refinement of methodological plans and structuring the multi round search for improved procedures. Section \ref{sec:optim_method} presents the optimization strategies that implement this capability. Scientific tools may be invoked within this subsystem when computational or empirical assessment is required.

\textbf{\textit{The third capability is long horizon memory.}} It supports the Evolution subsystem by maintaining persistent storage and retrieval of contextual information, reasoning traces, and experimental outcomes. Section \ref{sec:memory_method} describes its structured organization and interaction rules.

Across these capabilities, {\ProjectName} maintains the continuity, adaptability, and scalability needed for reliable and continuously improving scientific discovery.


%%%%%%%%%%%%%%%%%%%%%%%%%%%%%%%%%%%%%%%%%%%%%%%%%%%%%%%%%%%%%%%%%%%%%%%%%%%%%%%%
% \subsection{Cross Disciplinary Deep Search and Graph Construction \\
% \normalsize \color{gray} \textbf{\textit{Deep Research Capability within the Generation Subsystem}}}
\subsection{Cross Disciplinary Graph Construction and Knowledge Capturing}
\vspace{-0.5em}
{\raggedleft \normalsize \color{gray} \textbf{\textit{Deep Research Capability within the Generation Subsystem}}\par}


\label{sec:dr_method}
To enable cross disciplinary knowledge construction and utilization, our design operates on both data and methodological levels. On the data side, the system integrates diverse scientific sources with the assistance of domain specific tools to parse, normalize, and structure scientific information into a large scale multidisciplinary knowledge graph. On the methodological side, it identifies relations and dependencies across domains through a structured extraction workflow that combines model driven analysis with tool assisted processing of specialized scientific data, enabling deep and effective cross disciplinary knowledge integration.

\subsubsection{Cross-Disciplinary Knowledge Graph}
To support accurate and comprehensive deep research, we maintain a cross-disciplinary knowledge graph (KG). Notably, it differs from traditional KGs \cite{DBLP:conf/acl/ChamiWJSRR20,DBLP:conf/kdd/CaoLWL24}, which represent knowledge as triples including simple entities and relations; our KG instead captures a richer set of scientific elements. 

\paragraph{Graph Construction} 
From parsed outputs of papers, survey articles, technical reports, and domain notes, we construct a heterogeneous graph with nodes representing documents, key concepts, methods, datasets, empirical settings, and problem statements. The parsing process incorporates domain specific scientific tools to assist in identifying specialized entities and scientific attributes that are difficult to extract through general text analysis alone. Edges encode typed relations such as “cites’’ and “by product.’’ This design lets a single research idea sit at the intersection of multiple methodological and application communities and converts a flat corpus into a structured map where cross field dependencies emerge as paths rather than isolated points. For example, given the text
\textit{
``Esterification generally refers to the reaction of an alcohol with a carboxylic acid to give an ester and water. Common fats are esters; they can be hydrolysed back to alcohol and acid. A typical fat is triacylglycerol, formed from glycerol (propane 1,2,3 triol) and fatty acids (alkanoic acids containing 4–28 carbon atoms).’’}
the corresponding structured representation is shown in Fig.~\ref{img:kg-illustration}.

\begin{figure}[t]
\vspace{-2pt}
\begin{center}
% \includegraphics[width=\textwidth]{figures/performance.pdf} 
\includegraphics[width=\textwidth]{figures/KG_fig.pdf} 
\end{center}
\vspace{-10pt}
\caption{The illustration for our cross-disciplinary knowledge graph.}\label{img:kg-illustration}
\vspace{-4pt}
\end{figure}

\begin{figure}[h]
\vspace{-2pt}
\begin{center}
% \includegraphics[width=\textwidth]{figures/performance.pdf} 
\includegraphics[width=\textwidth]{figures/internagent-dr-case-study_cropped.pdf} 
\end{center}
\vspace{-10pt}
\caption{The illustration for our flow graph.}\label{img:flow_graph}
\vspace{-4pt}
\end{figure}


\paragraph{Knowledge Extraction and Retrieval} We employ a schema-guided extraction workflow to build a knowledge graph from noisy, cross-domain text. First, candidate entities are identified via domain-agnostic named entity recognition and noun-phrase mining, and document-level co-citation and co-usage relations are used to establish initial concept links. A subsequent consolidation step refines node types and edge semantics, aligning textual evidence with citation evidence into a compact cross-disciplinary graph. To answer deep research queries, we integrate graph search with dense vector retrieval: graph search uncovers the nodes and paths that connect the query to relevant methods and domains, while dense retrieval captures semantically related items not directly linked in the graph. Finally, a ranking step merges these results and outputs path-structured evidence chains, which the deep research module then analyzes to reveal cross-disciplinary connections.

\subsubsection{Flow Graph}
In real scientific deep research tasks, knowledge often exhibits highly non-linear and dynamic dependencies. Conventional sequential research process struggle to capture these relationships effectively, which can lead to redundant information, over-reliance on early hypotheses, and inflexible reasoning processes. To address these challenges, we introduce Dynamic Structured Knowledge Flow as a core principle of deep research system, enabling systematic and adaptive organization of knowledge throughout the research process. Specifically, we capture the knowledge in research process as a directed acyclic graph (DAG) that explicitly represents tasks, subtasks, and their dependencies. 

\paragraph{Structured Knowledge Flow}
The research process is organized as a directed graph, which provides a structured
representation of the reasoning process. Formally, the research process is
defined as:
\begin{equation}
G = (V, E),
\end{equation}
where $V$ denotes the set of nodes and $E$ denotes the set of directed edges encoding
dependencies among nodes. Each node $v_i \in V$ corresponds to a subtask or a key conceptual unit arising during the reasoning process. To explicitly capture its functional role and execution status, each node is represented as a tuple:
\begin{equation}
v_i = (t_i, d_i, s_i, c_i),
\end{equation}
where $t_i \in \{\text{search}, \text{solve}, \text{answer}\}$ specifies the task type
associated with the node, $d_i$ describes the task content, $s_i$ tracks the execution
state of the node, and $c_i$ stores the resulting knowledge context upon successful
completion of the task. Directed edges in the graph encode structural dependencies or reasoning constraints between nodes. Specifically, each edge is defined as:
\begin{equation}
e_{ij} = (v_i, v_j, r_{ij}) \in E,
\end{equation}
where $e_{ij}$ indicates a directed relationship from node $v_i$ to node $v_j$, and $r_{ij}$
characterizes the type of dependency between the two nodes, such as \emph{requires
result from}, \emph{provides evidence for}, or \emph{constrains reasoning of}. These
relational edges ensure that information and intermediate results are propagated in a
dependency-aware manner throughout the reasoning graph.


\paragraph{Dynamic planning and refinement}
The knowledge flow is constructed incrementally: starting from a root query node, a planner identifies nodes that require further decomposition or context enrichment, generates successor nodes, and updates dependency edges accordingly. This iterative expansion continues until the flow sufficiently covers all subproblems necessary to address the research objective. This design not only enables efficient multi-agent collaboration but also supports adaptive refinement of the knowledge flow as new evidence emerges, ensuring coherent, systematic, and verifiable reasoning throughout the research process.

\subsubsection{Graph-Guided Output Synthesis}
Building upon the dynamically constructed knowledge flow, we describes how the abstract graph structure is instantiated into concrete research outputs. Once the planner completes the incremental construction of the flow graph, executable nodes are activated according to their dependency states and progressively populated with domain knowledge and intermediate reasoning results. Through iterative node execution, state updating, and context propagation along graph edges, the system refines the structured knowledge flow and synthesizes the final research answer.

\paragraph{Cross-Disciplinary Knowledge Collector.}
To facilitate cross disciplinary insight, the Knowledge Collector gathers information from a diverse set of sources across multiple domains. These sources include outputs obtained through scientific tools and remote resources accessed via the Science Context Protocol (SCP)~\citep{jiang2025scp}. By integrating multi domain knowledge, the system can uncover unexpected connections and inspire ideas that may not emerge within a single discipline. Executable nodes with satisfied dependencies are assigned to agents, which decompose each subtask into a sequential reasoning and information retrieval process, iteratively assembling the knowledge required to resolve it. After a node is executed, its state is updated and the resulting knowledge context is propagated to dependent nodes, ensuring that subsequent reasoning benefits from the most up to date and contextually rich information. This design enables structured, adaptive, and collaborative knowledge acquisition throughout the research process.

\paragraph{Reasoning capability enhancement}
We adopt a reasoning capability enhancement strategy that enables reasoning along multiple complementary pathways. For a given query, the model generates three forms of responses: a direct answer based solely on the input, a search augmented answer that incorporates evidence retrieved from external sources and scientific tools, and a self driven answer obtained through internal retrieval and refinement. These complementary outputs are aggregated to form the final response, balancing intrinsic reasoning, external evidence, and self consistency. This multi path strategy improves answer completeness and factual reliability while reducing reliance on any single reasoning pathway.


%%%%%%%%%%%%%%%%%%%%%%%%%%%%%%%%%%%%%%%%%%%%%%%%%%%%%%%%%%%%%%%%%%%%%%%%%%%%%%%%
% \subsection{Experiment Execution and Multi-round Parallel Optimization \\
% \normalsize \color{gray} \textbf{\textit{solution refinement Capability within the Verification Subsystem}}}
\subsection{Experiment Execution and Multi-round Parallel Optimization}
\vspace{-0.5em}
{\raggedleft \normalsize \color{gray} \textbf{\textit{Solution Refinement Capability within the Verification Subsystem}}\par}

\label{sec:optim_method}

% \textbf{[UNDER REORGANIZATION]: MLE-bench removed}
The transformation of a refined methodology into a verifiable scientific result requires an efficient and reliable validation loop. In both computational algorithm design and physical wet‑lab experimentation, the search space of possible configurations is extremely large, and linear trial‑and‑error procedures often converge prematurely. This section introduces the multi‑round parallel experiment optimization and execution framework, which enables {\ProjectName} to explore this space autonomously and progressively converge toward high‑quality experimental proposals.

\subsubsection{Generative Design for Experimental Optimization}

Efficiently exploring a large and unstructured design space is a central challenge in automated scientific experimentation. Traditional strategies \cite{aider, automind} based on linear refinement or tree‑structured search often face three fundamental limitations: %~\cite{xx}~\cite{xx}
\textbf{Isolated Trajectories} arise when insights discovered in one search path cannot inform parallel explorations.
\textbf{Unexploited Search History} occurs when informative patterns across longer trajectories are not captured or reused.
\textbf{Limited Idea Composition} restricts the integration of promising elements from different branches into improved solutions.

We formalize the experimental optimization problem as identifying the optimal solution within a search space $\mathcal{S}$, where each candidate solution $s \in \mathcal{S}$ represents a complete experimental proposal, including code logic, parameter configurations, and physical operation protocols. The objective is to find:
\begin{equation}
s^* = \arg\max_{s \in \mathcal{S}} h(\mathcal{T}, s),
\label{eq:optimization_objective}
\end{equation}
where $h(\mathcal{T}, s)$ denotes the evaluation metric of solution $s$ on a given task $\mathcal{T}$.

To address the limitations of conventional search, {\ProjectName} adopts a \textit{Graph‑Augmented Monte Carlo Search} framework. This approach preserves the exploration–exploitation balance of Monte Carlo Tree Search while replacing its rigid tree structure with a dynamic solution graph that aggregates information across all prior experiments. The search still follows the classical loop of selection, expansion, simulation, and backpropagation, but its effectiveness comes from a strengthened expansion phase powered by four graph‑based operators:
\begin{itemize}
    \item \textbf{Primary Expansion.} Generates a new proposal using only its immediate parent. It performs localized adjustments such as parameter refinement or correction of logical inconsistencies, creating the core backbone of parent and child edges used in credit assignment.
    
    \item \textbf{Intra‑branch Evolution.} Conditions proposal generation on both the parent and the historical trajectory of ancestors within the same branch. By analyzing recent successes and failures, it reinforces productive design changes and avoids repeatedly exploring unpromising strategies, formalizing a localized form of self‑reflection.
    
    \item \textbf{Cross‑branch Reference.} Introduces targeted transfer of design elements across different branches. When a branch stagnates, the system identifies high‑performing nodes elsewhere in the solution graph and uses them as references, allowing the new proposal to incorporate robust structural patterns or methodological modules discovered in parallel explorations.
    
    \item \textbf{Multi‑branch Aggregation.} Synthesizes complementary strengths from multiple top‑performing nodes across the solution graph. By decomposing these proposals into their essential components and recombining them, the operator produces hybrid designs that integrate successful ideas from previously independent search trajectories.
\end{itemize}

Once a new proposal is generated through one of these operators, it is executed in the corresponding environment, either a computational simulator or a physical experimental system, to obtain an evaluation score. This score is backpropagated through the proposal’s ancestral path to guide subsequent exploration. By integrating graph‑based information flow into the Monte Carlo search process, {\ProjectName} transforms experimental design into a collaborative and cumulative optimization pipeline, enabling rapid convergence toward high‑quality scientific solutions.

\subsubsection{Scenario}
\noindent \textbf{Code Opimization for Algorithm Discovery}
In algorithm discovery tasks, each proposal is an executable program specifying data‑processing steps, model components, and evaluation settings. The search module generates new variants by refining computational logic or integrating effective structures identified in other branches. Each candidate is executed in a controlled environment that compiles the code and evaluates its performance on benchmark datasets. Quantitative metrics such as accuracy, runtime, and resource usage are returned to the optimization module and backpropagated through the proposal’s lineage, enabling systematic refinement of algorithmic designs.

\noindent \textbf{Experimental Optimization for Empirical Discovery}
In empirical discovery tasks, each proposal specifies a full experimental procedure that may be executed either through computational simulation or on physical laboratory systems. The search process refines these procedures by adjusting parameters, modifying operational steps, or incorporating effective sub‑protocols identified across branches. When a proposal is simulated, domain models predict experimental outcomes such as reaction yield or protein stability. When it is executed physically, SCP \cite{jiang2025scp} coordinates automated devices to perform the protocol and collect measurements such as fluorescence intensity or assay signal quality. All quantitative results, whether simulated or physically measured, are returned to the optimization module for backpropagation, enabling iterative improvement of empirical workflows.


%%%%%%%%%%%%%%%%%%%%%%%%%%%%%%%%%%%%%%%%%%%%%%%%%%%%%%%%%%%%%%%%%%%%%%%%%%%%%%%%
% \subsection{Structured Cognitive Memory for Long Horizon Scientific Discovery \\
% \normalsize \color{gray} \textbf{\textit{Long Horizon Memory Capability within the Evolution Subsystem}}}
\subsection{Structured Cognitive Memory for Long Horizon Scientific Discovery}
\vspace{-0.5em}
{\raggedleft \normalsize \color{gray} \textit{\textbf{Long Horizon Memory Capability within the Evolution Subsystem}}\par}


\label{sec:memory_method}

\begin{figure}[h]
\vspace{-2pt}
\begin{center}
% \includegraphics[width=\textwidth]{figures/performance.pdf} 
\includegraphics[width=\textwidth]{figures/mem_structure.pdf} 
\end{center}
\vspace{-10pt}
\caption{The illustration for our Structured Cognitive Memory.}\label{img:mem_structure}
\vspace{-4pt}
\end{figure}

To support adaptive, efficient, and long horizon scientific discovery, {\ProjectName} incorporates a hierarchical memory subsystem referred to as Structured Cognitive Memory. This subsystem is engaged throughout the entire discovery loop and maintains continuity across cycles, allowing the agent to operate coherently over extended durations. It consolidates procedural, episodic, and semantic information into a unified structure that enables short term refinement, mid term adaptation, and long term conceptual development. The overall framework of Structured Cognitive Memory is shown in Fig.~\ref{img:mem_structure}.

\subsubsection{Strategy-Procedural Memory}
Strategy Procedural Memory (SPM) supports the deep research capability that InternAgent-1.5 invokes throughout the entire scientific discovery process whenever complex analytical reasoning is required. This capability involves integrating literature evidence, synthesizing contextual knowledge, constructing coherent multi-step reasoning plans, correcting failed strategies from earlier research workflows, and analyzing the root causes behind those failures. Instead of storing raw trajectories, the system distills reusable procedural structures from past reasoning processes, including both validated effective patterns and lessons learned from unsuccessful attempts. These procedural structures capture the key decision pivots, organizational patterns that have proven effective across earlier research workflows, as well as the diagnostic insights into failed strategies and their underlying reasons, serving as strategic priors that can be applied flexibly at different stages of the pipeline to avoid recurring pitfalls and optimize reasoning paths.

Given a historical trajectory $T$, SPM constructs a compact representation as follows:
\begin{equation}
\mathbf{z}_{T} = f_{\text{proc}}(T),
\end{equation}
which captures essential procedural states extracted from the full execution trace.  
When a new deep research query $q$ arrives, {\ProjectName} retrieves trajectories with procedurally aligned structures:
\begin{equation}
\mathcal{S}(q) =
\operatorname*{topk}_{T \in \mathcal{M}_{\text{SPM}}}
\mathrm{sim}\!\left( f_{\text{proc}}(q),\, \mathbf{z}_T \right).
\end{equation}
These retrieved strategic priors guide the planner toward globally coherent reasoning graphs, while constraining the executor to avoid redundant execution steps and unnecessary tool calls, thereby providing a stable and efficient procedural foundation for the downstream discovery process.

\subsubsection{Task-Episodic Memory}
Task Episodic Memory (TEM) supplies fine grained, within trajectory evidence that enables rapid adaptation during iterative experimentation. After each experiment, the system stores an episodic unit containing the attempted method $m$, extracted metrics $y$, and an improvement judgment. Each unit is encoded using a hybrid representation that combines semantic embeddings with sparse lexical features to support both conceptual and literal alignment.

During hypothesis refinement, relevant episodes are retrieved through the following formula:
\begin{equation}
\mathcal{R}(q) =
\operatorname*{topk}_{e \in \mathcal{E}}
\mathrm{sim}\!\left( f_{\text{enc}}(q),\, f_{\text{enc}}(e) \right),
\end{equation}
where $q$ denotes the current hypothesis.  
The retrieved episodes are injected directly into the generation context, helping the system avoid unsuccessful methodological directions, exploit successful ones, and refine hypotheses efficiently within each research trajectory.

\subsubsection{Semantic-Knowledge Memory}

Semantic Knowledge Memory (SKM) consolidates conceptual information across sessions and supports the long horizon evolution of research objectives. It consists of a Long Term Experience Library, which stores distilled methodological insights, and an Idea Graph that tracks the semantic topology of previously explored research directions.Specifically, upon the end of each experimental batch, the system employs a pairwise combination strategy for the generated methods to maximize the utilization of existing information. By leveraging contrastive learning between each combination according to their methods and experimental results, {\ProjectName} extracts both high-level methodological principles and low-level experimental heuristics to construct Long-term Experience Library.
Given a research goal $G$, long term knowledge entries are retrieved via the following formula:
\begin{equation}
\mathcal{K}(G) =
\operatorname*{topk}_{k \in \mathcal{L}}
\mathrm{sim}\!\left( f_{\text{enc}}(G),\, f_{\text{enc}}(k) \right).
\end{equation}

To promote continued exploration, each candidate objective $c$ is assigned a novelty score:
\begin{equation}
\mathrm{nov}(c) =
1 - \max_{x \in \mathcal{G}}
\mathrm{sim}\!\left( f_{\text{enc}}(c),\, f_{\text{enc}}(x) \right),
\end{equation}
which encourages the selection of objectives that extend beyond previously explored conceptual regions.  
In this way, SKM provides the semantic continuity and innovation pressure required for sustained multi-session scientific discovery.
